\section{Background}
\label{sec:background}

\subsection{E-Scooter Safety and Speed-Related Injuries}
\label{sec:bg_safety}

The rapid expansion of shared e-scooters since 2017 has been accompanied by a well-documented surge in injuries. Hospital-based studies established that e-scooter injuries predominantly involve upper-extremity fractures and head trauma, with head injuries accounting for approximately 22\% of presentations \citep{trivedi2019injuries, badeau2019emergency, bloom2021standing}. Systematic reviews confirm these patterns: \citet{singh2022impact} found bony injuries as the most common type (39.2\%), while \citet{spota2024injury} identified two prevalent injury profiles---mild upper-limb injuries and severe head trauma. The economic burden is substantial, with \citet{kahan2024escooter} reporting annual e-scooter hospital charges of USD~10.4~million in a single US city.

Speed is implicated as a primary risk factor across this literature. The power model of the speed--crash relationship \citep{elvik2013reexamination} predicts that injury severity increases as the fourth power of impact speed---a relationship especially consequential for unprotected e-scooter riders. \citet{yang2020safety} found that fatal e-scooter crashes disproportionately involved high-speed collisions with motor vehicles. Yet despite the central role of speed in crash severity, empirical evidence on actual riding speeds and regulatory compliance remains surprisingly sparse.

\subsection{Speed Governance: Regulation Versus Technological Enforcement}
\label{sec:bg_governance}

Jurisdictions worldwide have imposed regulatory speed limits on e-scooters: 25~km/h in South Korea \citep{koreaherald2020escooter} and France, 20~km/h in Germany and the UK, and 15--25~km/h across US cities \citep{ma2021municipal, sexton2023shared}. These limits are typically enforced through two mechanisms: traffic regulations (penalties for exceeding the limit) and technological enforcement (software-based speed governors embedded in the vehicle firmware). In practice, traffic enforcement of e-scooter speed limits is rare---police resources are concentrated on motor vehicle enforcement, and speed cameras do not typically capture e-scooter license plates. This makes software governors the \textit{de facto} primary enforcement mechanism.

Despite their prevalence, the empirical evidence base for governor effectiveness is thin. \citet{cicchino2023speed} measured 2,004 riders' speeds via roadside observation in two US cities with different speed limiter settings, finding that governor settings meaningfully affect behavior. However, between-city comparisons conflate governor effects with infrastructure and demographic differences. Technology-based speed management, including geofencing \citep{lime2023geofencing, caggiani2023geofencing}, has received growing attention as a policy tool, but large-scale behavioral evaluations remain limited \citep{neshagaran2024safety}. No study has evaluated governor effectiveness using within-operator variation---where identical hardware operates under different software speed limits---or provided causal evidence via an exogenous policy change.

Beyond direct speed effects, the broader questions of behavioral compensation and demand responses to governance remain unresolved (Sections~\ref{sec:bg_compensation} and \ref{sec:bg_demand}). Speed variability, an independent risk indicator for micromobility \citep{twisk2021speed}, could increase even if maximum speeds are constrained, and governance could reduce ridership if governed vehicles are perceived as less desirable.

\subsection{Risk Compensation and the Peltzman Hypothesis}
\label{sec:bg_compensation}

The Peltzman hypothesis \citep{peltzman1975effects} posits that mandated safety improvements may induce offsetting behavioral responses: because individuals adjust risk-taking in response to changes in perceived safety, interventions that reduce risk on one margin may encourage risk-taking on others. In his analysis of automobile safety regulation, Peltzman argued that while seatbelt mandates reduced fatality rates per crash, they also increased crash frequency---suggesting partial behavioral offset. \citet{wilde1982theory} formalized this intuition as ``risk homeostasis theory'': individuals maintain a target level of acceptable risk, and safety interventions that reduce risk below this target elicit compensatory behavior until the target is restored.

Empirical evidence on risk compensation is domain-dependent and contested. In the automotive context, the most comprehensive synthesis---\citeauthor{elvik2013reexamination}'s power model meta-analysis---finds that engineering interventions (road design, vehicle crashworthiness) produce safety benefits substantially larger than any compensatory offset. Seat belt wearing, the original Peltzman context, shows a more nuanced picture: \citet{cohen2003effects} found that US mandatory seatbelt laws reduced overall fatalities with no evidence of offsetting increases in crashes. \citet{hedlund2000risky} argued that behavioral adaptation is most plausible when the intervention is \textit{salient} (the user knows protection has increased) and the regulated behavior is \textit{voluntary}---conditions that may not apply to speed governors, which operate invisibly.

For e-scooter speed governance specifically, risk compensation has not been tested. The governor context offers a particularly clean test for three reasons: (1) the intervention is binary and precisely defined (speed cap present or absent), (2) multiple behavioral margins are observable simultaneously from GPS data---enabling a comprehensive multi-margin test rather than the single-outcome designs common in the compensation literature, and (3) the within-user mode-switcher design cleanly separates compositional shifts (different riders entering the pool) from genuine behavioral adaptation (same riders changing behavior). This separation is critical because aggregate-level analyses can confound composition with compensation, as we demonstrate in Section~\ref{sec:res_compensation}.

\subsection{Demand Responses to Safety Regulation}
\label{sec:bg_demand}

A distinct but related concern is whether safety regulation reduces participation in the activity it regulates---the ``demand cost'' of safety. This margin is analytically separate from risk compensation: compensation concerns behavioral changes \textit{within} the activity (riding more aggressively on governed vehicles), while demand responses concern changes in \textit{whether} individuals participate at all.

The demand concern is empirically grounded in other transportation domains. Mandatory bicycle helmet laws in Australia were associated with 20--35\% reductions in cycling participation \citep{robinson1996head}, raising concerns that the health benefits of helmets were partially offset by reduced physical activity. For motorcycle helmet laws, \citet{dee2009motorcycle} found that universal helmet mandates reduced motorcyclist fatalities but also reduced registered motorcycles---suggesting a participation margin response. Graduated driver licensing (GDL) programs, which restrict novice drivers' exposure to high-risk conditions, have been shown to reduce teen crash involvement without commensurate increases in crashes after restrictions are lifted \citep{shope2007graduated}.

In the e-scooter context, the demand question is whether speed governance reduces ridership. If governed e-scooters are perceived as slower, less enjoyable, or less useful for commuting, riders may reduce trip frequency, shorten trips, or abandon the platform entirely. This concern is particularly salient for shared micromobility operators, whose revenue depends on trip volume. Yet no empirical evidence exists on the demand elasticity of e-scooter speed governance. Our multi-outcome framework directly addresses this gap by including trip counts and active user counts as causal outcomes alongside safety and dynamics margins.

\subsection{Natural Experiments in Transportation Safety}
\label{sec:bg_natural}

Natural experiments---where an exogenous event creates treatment and control conditions without researcher inter\-vention---have a long history in transportation safety. Speed limit changes on highways, helmet legislation, seatbelt mandates, and alcohol regulation have all been evaluated using before--after or difference-in-differences (DiD) designs \citep{elvik2013reexamination}. In the e-scooter domain, however, natural experiments are rare because policy changes typically affect all operators simultaneously within a jurisdiction, eliminating within-jurisdiction variation for causal identification.

The December 2023 TUB ban by Swing presents an unusual opportunity. The policy change was system-wide (all 52 cities simultaneously), but the treatment intensity varied across cities because pre-ban TUB adoption rates differed substantially (32--70\% of trips). This creates a continuous treatment variable---cities with higher pre-ban TUB share experienced larger ``doses'' of governor imposition---enabling a dose-response design within a standard two-way fixed effects DiD framework. The staggered-adoption concerns that complicate many DiD applications \citep{roth2022pretest} do not apply here because the treatment date is identical across all cities.

\subsection{Research Gaps}
\label{sec:bg_gaps}

Our review identifies three interconnected gaps. First, no prior study has analyzed e-scooter speed behavior at a scale sufficient to characterize population-level patterns---existing work relies on spot-speed observations ($n < 2{,}000$) or small naturalistic samples ($n < 500$ trips). Second, no \textit{causal} evidence exists for speed governor effectiveness; all prior comparisons are cross-sectional. Third, behavioral compensation after governor policy changes has not been tested. We address all three gaps using 21~million trips with continuous speed profiles, a natural experiment with a continuous treatment design, and multi-outcome DiD with within-user mode-switcher analysis. Table~\ref{tab:lit_comparison} positions our study relative to prior work.

\begin{table}[htbp]
\centering
\caption{Comparison with closely related studies on e-scooter speed behavior.}
\label{tab:lit_comparison}
\small
\begin{tabular}{lrrcccc}
\toprule
\textbf{Study} & \textbf{N} & \textbf{Cities} & \textbf{Speed data} & \textbf{Governor} & \textbf{Causal} & \textbf{Compensation} \\
\midrule
\citet{cicchino2023speed} & 2,004 & 2 & Spot & Between-city & No & No \\
\citet{ma2021escooter} & $<$500 & 1 & Profile & No & No & No \\
\citet{jiao2020shared} & 900K & 1 & None & No & No & No \\
\citet{mckenzie2019spatiotemporal} & 137K & 1 & None & No & No & No \\
\citet{neshagaran2024safety} & 2,400 & 2 & Observed & No & No & No \\
\addlinespace
\textbf{This study} & \textbf{21M} & \textbf{52} & \textbf{Profile} & \textbf{Within-operator} & \textbf{DiD} & \textbf{Yes} \\
\bottomrule
\end{tabular}
\end{table}
